Una empresa de transporte llamada \textbf{CargoPlan} se encarga de distribuir productos en tres zonas diferentes: Zona A, Zona B y Zona C. 
Cada zona tiene una demanda específica y el transporte requiere recursos limitados en términos de combustible y horas de conducción.

Los límites semanales para estos recursos son \textbf{8000 litros de combustible y 3000 horas de conducción}. El beneficio por tonelada transportada a cada zona es de \textbf{50, 40 y 60 euros} para las zonas A, B y C, respectivamente.

Cada tonelada transportada consume los siguientes recursos. Zona A: \textbf{6} litros de combustible y \textbf{3} horas de conducción. Zona B: \textbf{4} litros de combustible y \textbf{2} horas de conducción. Zona C: \textbf{8} litros de combustible y \textbf{5} horas de conducción.

Además, existe un contrato que obliga a transportar al menos \textbf{500 toneladas en total} a las zonas B y C combinadas.


Si $x_1$, $x_2$ y $x_3$ representan las toneladas transportadas semanalmente a las zonas A, B y C, respectivamente, el modelo de programación lineal que permite maximizar el beneficio total es el siguiente:

\begin{equation}
\begin{split}
\mbox{max. } z = 50x_1 + 40x_2 + 60x_3\\
s.a.:\\
6x_1 + 4x_2 + 8x_3 \leq 8000\\
3x_1 + 2x_2 + 5x_3 \leq 3000\\
        x_2 +  x_3 \geq 500\\
x_1,\,\,x_2,\,\,x_3\geq 0\\
\end{split}
\end{equation}

Se dispone de la siguiente tabla, correspondiente a una solución obtenida mediante la aplicación del método del simplex.

\begin{center}
\begin{tabular}{c|c|ccccccc|}
 & $z$ & $x_1$ & $x_2$ & $x_3$ & $h_1$ & $h_2$ & $h_3$ & $a_3$\\ 
 & -60000 & -10 & 0 & -40 & 0 & -20 & 0 & 0 \\ 
\hline
$h_1$ & 2000  & 0 & 0 & -2 & 1 & -2 & 0 & 0\\ 
$h_3$ & 1000  & $3/2$ & 0 & $3/2$ & 0 & $1/2$ & 1 & -1\\ 
$x_2$ & 1500  & $3/2$ & 1 & $5/2$ & 0 & $1/2$ & 0 & 0\\ 
\hline
\end{tabular}
\end{center}





Se pide:

\begin{enumerate}
    \item Interpretar la solución óptima obtenida, incluyendo el beneficio total, la asignación de toneladas transportadas a cada zona y el uso de los recursos disponibles.
\end{enumerate}

Cada una de las preguntas siguientes se refieren a la solución óptima del problema.

\begin{enumerate}
    \setcounter{enumi}{1} % Ajusta el contador para iniciar en 2

    \item Identificar dentro de qué rango podría variar la disponibilidad de combustible (en litros) sin que se modifiquen las variables básicas de la solución óptima y explicar qué cambiaría en la solución el cambio de la disponibilidad de combustible dentro de dicho rango. ¿Cuál sería el plan de transporte fuera de ese rango?

    \item Se ha hecho un análisis del tiempo de conducción en la zona A para reducir el tiempo actual de 3 horas por cada tonelada transportada. ¿Por debajo de qué valor de tiempo de conducción en la zona A sería interesante transportar sus productos a la zona A?

    \item Cuánto estaría dispuesta a pagar la empresa por disponer de una hora adicional de conducción cada semana.

  % \item Cargoplan está renegociando los precios de los clientes de la zona B. Dentro de qué rango de valores del ingreso por tonelada en la zona B se mantendría el plan de distribución.

    \item Evaluar el impacto de relajar o de hacer más restrictivo el contrato que exige transportar al menos 500 toneladas en total a las zonas B y C. En particular, indicar cuánto se podría relajar o cuánto más restrictivo podría ser el acuerdo sin que cambie el plan de distribución óptimo.

    \item Obtener el nuevo plan óptimo de transporte si, por razones logísticas, la empresa decidiera transportar al menos tantas toneladas a la zona A como a la zona C.

    % \item Obtener el nuevo plan óptimo de transporte si, por razones logísticas, la empresa decidiera transportar las mismas cantidades a la zona A que a la zona B.
\end{enumerate}
